% GENERATED FILE. Edit canonical lessons, then run npm run generate:books.
% canonical-source-hash: fnv1a64:affabb6909a48b6e
% canonical-lessons: ML-C479-alayuka, ML-C479-tannuka, ML-C479-sahayatrikan, ML-C479-itattavalam, ML-C479-pinnituka

\chapter{On the Road for Days}
\label{ch:ml-on-the-road-for-days}
\hlchaptermodality{479}

\begin{chapteropening}
\textbf{By the end of this chapter:} \emph{I can say to wander, to stay over, a fellow traveller, a stopover and to pass by.}

Complete the last lesson of chapter 479: I can say to wander, to stay over, a fellow traveller, a stopover and to pass by.
\end{chapteropening}

\section[\texorpdfstring{\ml{അലയുക} (alayuka)}{അലയുക (alayuka)}]{\ml{അലയുക} (alayuka) — to wander; to roam}

\label{lesson:ML-C479-alayuka}



\begin{quote}
Before the new one: say the Malayalam for an itinerary, then the Malayalam for to roam around.
\end{quote}

\subsection*{You'll want to know: \ml{അലയുക}}
\textbf{\ml{അലയുക}} — \emph{alayuka} — \textquotedblleft{}to wander; to roam\textquotedblright{}.

\begin{quote}
\ml{വഴി} \ml{അറിയാതെ} \ml{ഞങ്ങൾ} \ml{ഒരു} \ml{മണിക്കൂർ} \ml{അലഞ്ഞു}.
\end{quote}

\begin{quote}
\emph{vaḻi aṟiyāte ñaṅṅaḷ oru maṇikkūr alaññu.}
\end{quote}

\begin{quote}
Not knowing the way, we wandered for an hour.
\end{quote}

Past: \ml{അലഞ്ഞു}, \emph{alaññu}.

It means wandering without a plan, often tired or lost.

\begin{cousinweb}
\textbf{\ml{അലയുക}} (\emph{alayuka}) is native Dravidian; DEDR gives the Malayalam verb as to be tossed about and to roam (DEDR 240).
\end{cousinweb}

\begin{grammarlens}[title={What you've built}]
One more word for getting around.
\end{grammarlens}

\subsection*{Your turn}
\begin{itemize}
\raggedright
  \item \emph{Say it:} \emph{alayuka}
  \item \emph{Say it:} \emph{alayuka}, once more
  \item \emph{Say it:} \emph{vaḻi aṟiyāte ñaṅṅaḷ oru maṇikkūr alaññu.} — Not knowing the way, we wandered for an hour.
  \item \emph{Recall:} say the Malayalam for an itinerary, then the Malayalam for to roam around, then say \emph{alayuka} again
\end{itemize}

\begin{tcolorbox}[breakable,colback=teal!4,colframe=teal!35!black,title={Before you move on}]
What does \ml{അലയുക} mean? (\textquotedblleft{}To wander; to roam\textquotedblright{}.) Say it once more.
\end{tcolorbox}
\section[\texorpdfstring{\ml{തങ്ങുക} (taṅṅuka)}{തങ്ങുക (taṅṅuka)}]{\ml{തങ്ങുക} (taṅṅuka) — to stay over; to halt}

\label{lesson:ML-C479-tannuka}



\begin{quote}
Before the new one: say the Malayalam for to roam around, then the Malayalam for to wander.
\end{quote}

\subsection*{You'll want to know: \ml{തങ്ങുക}}
\textbf{\ml{തങ്ങുക}} — \emph{taṅṅuka} — \textquotedblleft{}to stay over; to halt\textquotedblright{}.

\begin{quote}
\ml{യാത്രയ്ക്കിടെ} \ml{ഞങ്ങൾ} \ml{ഒരു} \ml{രാത്രി} \ml{ബന്ധുവിന്റെ} \ml{വീട്ടിൽ} \ml{തങ്ങി}.
\end{quote}

\begin{quote}
\emph{yātraykkiṭe ñaṅṅaḷ oru rātri bandhuvinṟe vīṭṭil taṅṅi.}
\end{quote}

\begin{quote}
On the journey we stayed a night at a relative's house.
\end{quote}

Past: \ml{തങ്ങി}, \emph{taṅṅi}.

\ml{താമസിക്കുക} is to live somewhere; \ml{തങ്ങുക} is to stop for a short while on the way.

\begin{cousinweb}
\textbf{\ml{തങ്ങുക}} (\emph{taṅṅuka}) is native Dravidian, \textquotedblleft{}to stay, to stop, to rest\textquotedblright{} (DEDR 3014).
\end{cousinweb}

\begin{grammarlens}[title={What you've built}]
One more word for getting around.
\end{grammarlens}

\subsection*{Your turn}
\begin{itemize}
\raggedright
  \item \emph{Say it:} \emph{taṅṅuka}
  \item \emph{Say it:} \emph{taṅṅuka}, once more
  \item \emph{Say it:} \emph{yātraykkiṭe ñaṅṅaḷ oru rātri bandhuvinṟe vīṭṭil taṅṅi.} — On the journey we stayed a night at a relative's house.
  \item \emph{Recall:} say the Malayalam for to roam around, then the Malayalam for to wander, then say \emph{taṅṅuka} again
\end{itemize}

\begin{tcolorbox}[breakable,colback=teal!4,colframe=teal!35!black,title={Before you move on}]
What does \ml{തങ്ങുക} mean? (\textquotedblleft{}To stay over; to halt\textquotedblright{}.) Say it once more.
\end{tcolorbox}
\section[\texorpdfstring{\ml{സഹയാത്രികൻ} (sahayātrikan)}{സഹയാത്രികൻ (sahayātrikan)}]{\ml{സഹയാത്രികൻ} (sahayātrikan) — a fellow traveller}

\label{lesson:ML-C479-sahayatrikan}



\begin{quote}
Before the new one: say the Malayalam for to wander, then the Malayalam for to stay over.
\end{quote}

\subsection*{You'll want to know: \ml{സഹയാത്രികൻ}}
\textbf{\ml{സഹയാത്രികൻ}} — \emph{sahayātrikan} — \textquotedblleft{}a fellow traveller\textquotedblright{}.

\begin{quote}
\ml{തീവണ്ടിയിലെ} \ml{സഹയാത്രികൻ} \ml{എനിക്ക്} \ml{ഒരു} \ml{പഴം} \ml{തന്നു}.
\end{quote}

\begin{quote}
\emph{tīvaṇṭiyile sahayātrikan enikkŭ oru paḻaṁ tannu.}
\end{quote}

\begin{quote}
My fellow traveller on the train gave me a banana.
\end{quote}

Plural: \ml{സഹയാത്രികർ}, \emph{sahayātrikar}.

\ml{സഹ}- means with, as in \ml{സഹപാഠി}, a classmate.

\begin{cousinweb}
\textbf{\ml{സഹയാത്രികൻ}} (\emph{sahayātrikan}) is built from Sanskrit \emph{saha}, \textquotedblleft{}with\textquotedblright{}, and \emph{yātrika}, \textquotedblleft{}relating to a journey\textquotedblright{} (Monier-Williams), with the masculine ending -\ml{ൻ}.
\end{cousinweb}

\begin{grammarlens}[title={What you've built}]
One more word for getting around.
\end{grammarlens}

\subsection*{Your turn}
\begin{itemize}
\raggedright
  \item \emph{Say it:} \emph{sahayātrikan}
  \item \emph{Say it:} \emph{sahayātrikan}, once more
  \item \emph{Say it:} \emph{tīvaṇṭiyile sahayātrikan enikkŭ oru paḻaṁ tannu.} — My fellow traveller on the train gave me a banana.
  \item \emph{Recall:} say the Malayalam for to wander, then the Malayalam for to stay over, then say \emph{sahayātrikan} again
\end{itemize}

\begin{tcolorbox}[breakable,colback=teal!4,colframe=teal!35!black,title={Before you move on}]
What does \ml{സഹയാത്രികൻ} mean? (\textquotedblleft{}A fellow traveller\textquotedblright{}.) Say it once more.
\end{tcolorbox}
\section[\texorpdfstring{\ml{ഇടത്താവളം} (iṭattāvaḷaṁ)}{ഇടത്താവളം (iṭattāvaḷaṁ)}]{\ml{ഇടത്താവളം} (iṭattāvaḷaṁ) — a stopover; a halting place}

\label{lesson:ML-C479-itattavalam}



\begin{quote}
Before the new one: say the Malayalam for to stay over, then the Malayalam for a fellow traveller.
\end{quote}

\subsection*{You'll want to know: \ml{ഇടത്താവളം}}
\textbf{\ml{ഇടത്താവളം}} — \emph{iṭattāvaḷaṁ} — \textquotedblleft{}a stopover; a halting place\textquotedblright{}.

\begin{quote}
\ml{ആ} \ml{ചെറിയ} \ml{ഗ്രാമം} \ml{ഞങ്ങളുടെ} \ml{യാത്രയിലെ} \ml{ഇടത്താവളമായിരുന്നു}.
\end{quote}

\begin{quote}
\emph{ā ceṟiya grāmaṁ ñaṅṅaḷuṭe yātrayile iṭattāvaḷamāyirunnu.}
\end{quote}

\begin{quote}
That small village was the stopover on our journey.
\end{quote}

\ml{ഇട} means the middle, as in \ml{ഇടയിൽ}, between; \ml{താവളം} is a camp or a place to stay.

\begin{cousinweb}
\textbf{\ml{ഇടത്താവളം}} (\emph{iṭattāvaḷaṁ}) is a Malayalam compound built from the native \ml{ഇട}, the middle or interval (DEDR 448), and \ml{താവളം}, a camp or lodging; the origin of \ml{താവളം} is not settled.
\end{cousinweb}

\begin{grammarlens}[title={What you've built}]
One more word for getting around.
\end{grammarlens}

\subsection*{Your turn}
\begin{itemize}
\raggedright
  \item \emph{Say it:} \emph{iṭattāvaḷaṁ}
  \item \emph{Say it:} \emph{iṭattāvaḷaṁ}, once more
  \item \emph{Say it:} \emph{ā ceṟiya grāmaṁ ñaṅṅaḷuṭe yātrayile iṭattāvaḷamāyirunnu.} — That small village was the stopover on our journey.
  \item \emph{Recall:} say the Malayalam for to stay over, then the Malayalam for a fellow traveller, then say \emph{iṭattāvaḷaṁ} again
\end{itemize}

\begin{tcolorbox}[breakable,colback=teal!4,colframe=teal!35!black,title={Before you move on}]
What does \ml{ഇടത്താവളം} mean? (\textquotedblleft{}A stopover; a halting place\textquotedblright{}.) Say it once more.
\end{tcolorbox}
\section[\texorpdfstring{\ml{പിന്നിടുക} (pinniṭuka)}{പിന്നിടുക (pinniṭuka)}]{\ml{പിന്നിടുക} (pinniṭuka) — to pass by; to leave behind}

\label{lesson:ML-C479-pinnituka}



\begin{quote}
Before the new one: say the Malayalam for a fellow traveller, then the Malayalam for a stopover.
\end{quote}

\subsection*{You'll want to know: \ml{പിന്നിടുക}}
\textbf{\ml{പിന്നിടുക}} — \emph{pinniṭuka} — \textquotedblleft{}to pass by; to leave behind\textquotedblright{}.

\begin{quote}
\ml{തീവണ്ടി} \ml{രണ്ട്} \ml{സ്റ്റേഷനുകൾ} \ml{പിന്നിട്ടു}.
\end{quote}

\begin{quote}
\emph{tīvaṇṭi raṇṭŭ sṟṟēṣanukaḷ pinniṭṭu.}
\end{quote}

\begin{quote}
The train passed two stations.
\end{quote}

Past: \ml{പിന്നിട്ടു}, \emph{pinniṭṭu}.

Literally to put behind: \ml{പിൻ}, behind, plus \ml{ഇടുക}, to put. A distance or a time can be passed too.

\begin{cousinweb}
\textbf{\ml{പിന്നിടുക}} (\emph{pinniṭuka}) is a native Dravidian compound of \ml{പിൻ}, behind (DEDR 4205), and \ml{ഇടുക}, to put (DEDR 442).
\end{cousinweb}

\begin{grammarlens}[title={What you've built}]
One more word for getting around.
\end{grammarlens}

\subsection*{Your turn}
\begin{itemize}
\raggedright
  \item \emph{Say it:} \emph{pinniṭuka}
  \item \emph{Say it:} \emph{pinniṭuka}, once more
  \item \emph{Say it:} \emph{tīvaṇṭi raṇṭŭ sṟṟēṣanukaḷ pinniṭṭu.} — The train passed two stations.
  \item \emph{Recall:} say the Malayalam for a fellow traveller, then the Malayalam for a stopover, then say \emph{pinniṭuka} again
\end{itemize}

\begin{tcolorbox}[breakable,colback=teal!4,colframe=teal!35!black,title={Before you move on}]
What does \ml{പിന്നിടുക} mean? (\textquotedblleft{}To pass by; to leave behind\textquotedblright{}.) Say it once more.
\end{tcolorbox}
